\chapter{Kết luận và hướng phát triển}
\section{Kết luận}
Đồ án đã hoàn thành việc thiết kế và thi công một hệ thống Gateway giám sát năng lượng điện dựa trên vi điều 
khiển ESP32-S3, đáp ứng đầy đủ các yêu cầu đề ra. Về phần cứng, thiết bị được thiết kế trên bo mạch PCB hai 
lớp tích hợp đầy đủ các khối chức năng: khối nguồn hạ áp ổn định, khối giao tiếp RS485 cách ly quang học, 
khối Ethernet sử dụng chip W5500, khối lưu trữ EEPROM và RTC, cùng giao diện người dùng LCD 16x04 với nút 
nhấn và encoder. Về phần mềm nhúng, hệ thống sử dụng FreeRTOS để quản lý đa nhiệm, cho phép thu thập dữ liệu 
định kỳ từ các đồng hồ đo điện qua giao thức Modbus RTU, lưu trữ dữ liệu vào thẻ nhớ SD Card khi mất kết nối 
mạng, và đồng bộ dữ liệu lên Cloud thông qua giao thức MQTT khi có kết nối. Ứng dụng cấu hình trên máy tính 
được xây dựng bằng PyQt cho phép kỹ thuật viên nhập các thông số thanh ghi và đồng bộ cấu hình xuống thiết bị 
qua Bluetooth BLE một cách thuận tiện. Trang web Dashboard phát triển trên nền React và Django hiển thị dữ 
liệu thời gian thực từ các đồng hồ điện, đồng thời cung cấp tính năng cảnh báo và mô phỏng trực quan trạng 
thái đường truyền Modbus. Nhìn chung, hệ thống vận hành ổn định và đáp ứng được nhu cầu thực tế của các doanh 
nghiệp trong việc giám sát và quản lý năng lượng điện theo thời gian thực.
\section{Hướng phát triển}
Trong tương lai, hệ thống có thể được cải tiến và mở rộng theo nhiều hướng để nâng cao tính ứng dụng và độ tin 
cậy. Thứ nhất, về phần cứng, có thể tích hợp thêm tính năng cập nhật firmware từ xa (OTA – Over The Air) giúp 
việc bảo trì và nâng cấp hệ thống trở nên thuận tiện hơn mà không cần tiếp cận trực tiếp thiết bị, đồng thời 
có thể mở rộng số lượng cổng RS485 để hỗ trợ thu thập dữ liệu từ nhiều đường dây Modbus song song. Thứ hai, 
về phần mềm và bảo mật, cần bổ sung cơ chế mã hóa dữ liệu truyền qua MQTT bằng giao thức TLS/SSL và xác thực 
người dùng trên trang web nhằm đảm bảo an toàn thông tin trong môi trường công nghiệp. Thứ ba, trang web 
Dashboard có thể được nâng cấp với các tính năng phân tích dữ liệu nâng cao như vẽ biểu đồ xu hướng tiêu thụ 
điện năng theo thời gian, tính toán chỉ số hiệu quả năng lượng (PUE) và xuất báo cáo tự động theo định kỳ. 
Cuối cùng, hệ thống có thể tích hợp thêm trí tuệ nhân tạo để phát hiện bất thường trong dữ liệu điện năng, 
từ đó cảnh báo sớm các sự cố tiềm ẩn trước khi chúng ảnh hưởng đến quy trình sản xuất của doanh nghiệp.